% ============================================================
% f2013_shaanxi_arts_q05_histogram.tex
% 2013 陕西卷（文史类）第 5 题 频率分布直方图（TikZ 矢量重绘）
% 分组 [10,15)、[15,20)、[20,25)、[25,30)、[30,35]
% 纵轴为 频率/组距；柱高依次为 0.02、0.04、0.06、0.05、0.03
% 横轴 0 与 10 之间有断裂记号
% ============================================================
\begin{tikzpicture}[
  x=1cm, y=1cm, line width=0.5pt, font=\footnotesize, >=Stealth
]

  % ---------- 坐标轴 ----------
  \draw[->] (0,0) -- (7.30,0);
  \draw[->] (0,0) -- (0,4.95);

  % 横轴原点附近的断裂记号
  \draw (0.16,0.09) -- (0.28,-0.09);
  \draw (0.34,0.09) -- (0.46,-0.09);

  % ---------- 频率/组距 纵轴标题 ----------
  \node at (-0.78,4.50) {$\dfrac{\text{频率}}{\text{组距}}$};

  % ---------- 五根柱 ----------
  \draw (0.80,0) rectangle (2.00,1.20);   % [10,15)
  \draw (2.00,0) rectangle (3.20,2.40);   % [15,20)
  \draw (3.20,0) rectangle (4.40,3.60);   % [20,25)
  \draw (4.40,0) rectangle (5.60,3.00);   % [25,30)
  \draw (5.60,0) rectangle (6.80,1.80);   % [30,35]

  % ---------- 纵轴刻度与虚线 ----------
  \foreach \y/\lab in {1.20/0.02, 1.80/0.03, 2.40/0.04, 3.60/0.06}{
    \draw (-0.07,\y) -- (0,\y);
    \node[left] at (-0.10,\y) {\lab};
    \draw[dashed, gray!55] (0,\y) -- (6.80,\y);
  }

  % ---------- 横轴刻度 ----------
  \foreach \x/\lab in {0/0, 0.80/10, 2.00/15, 3.20/20, 4.40/25, 5.60/30, 6.80/35}{
    \draw (\x,0) -- (\x,-0.08);
    \node[below] at (\x,-0.10) {\lab};
  }

  % ---------- 横轴标题 ----------
  \node[anchor=west] at (6.95,0.26) {长度/毫米};

\end{tikzpicture}
